\chapter{Parsing, bounds and unchanged output}
\MARSChapterHeading{Parsing, bounds and unchanged output}
\label{revision-lab:bounded-interfaces}

Import the development lesson \texttt{013-bounded-interfaces} and choose
Continue in Learn. Its complete \texttt{bounded-packet.asm} is also supplied in
\texttt{revision/labs/files}; both implement original project
\texttt{013-binary-packet-parser}, and the original checks still apply. Read the
main chapter's invented format and ordered decisions first. This supplement
explains the input, intermediate states and failure investigations alongside
the interactive route.

Learn follows an accepted packet through eight Debug stops, then a short header
and a bad checksum through two stops each. Its complete alternate sources are
\texttt{worked/packet-short.asm}, changing only available to seven, and
\texttt{worked/packet-checksum.asm}, changing only the checksum trailer. Keep
these files separate from the default program. Seven predictions and ordinary
Run of all three sources lead to independent binary delimiter parsing. When
working without Learn, make the corresponding copies and use the same observed
boundaries below.

\section{Predict bytes, fields and three lengths}

Write the default packet as twelve hexadecimal bytes and label header, payload
and trailer. Derive its declared payload length, consumed count, sequence number,
payload-sum residue and complete checksum residue. Distinguish the count of three
payload bytes from the count of twelve bytes needed to validate this packet.

Predict the result for available counts seven, eight, eleven and twelve without
changing the mapped packet bytes. Explain which reads have been justified at each
point. With eight bytes available, the declared body length can be decoded, but
the three payload bytes and trailer are not yet part of the supplied window.
A debugger can display those mapped bytes without granting the parser permission
to load them.

Repeat with a zero-length payload and a recomputed checksum. Eight bytes still
need more; nine can succeed. For the returned slice $input+8$, explain why a zero
payload length forbids treating its first addressed byte as payload. Then append
a second complete packet and predict the first call's consumed count. It should
not change merely because the caller supplied a longer window.

\section{Follow the first bound and the ordered errors}

For Learn's short-header route, open \texttt{worked/packet-short.asm} and
break at \texttt{parse\_packet}. The available argument is seven even though
the arena still contains the complete default packet. Step through the initial
comparison. The function must return need-more before
executing an input load. Stop at \texttt{parse\_ready}: \texttt{\$v0} is $-1$,
\texttt{\$v1} is zero, and all four output words retain their original values.
Continue through main and distinguish its status/pointer metadata stores from
forbidden writes inside the failed parser.

Figure~\ref{fig:revision-bounded-returns} compares the accepted twelve-byte
packet with a different incomplete prefix of eight bytes. Eight bytes establish
the whole header but still fail the complete-packet bound. The guided seven-byte
copy fails the earlier header bound before any input load. Both return
need-more, although they reach that result at different decisions.

\begin{figure}[htbp]
\centering
\begin{tikzpicture}[alt={Two actual light-theme Debug register observations at parse_ready. The valid twelve-byte packet returns v0 equal to 12 and v1 equal to 268501004, the input address plus eight. The same mapped bytes supplied with available count eight return v0 equal to minus one and v1 equal to zero. Native boxes identify the return registers.}]
\node[anchor=north west,inner sep=0,font=\small\bfseries] (atitle) at (0,0) {1. Twelve bytes supplied};
\node[anchor=north west,inner sep=0] (a) at ([yshift=-2mm]atitle.south west) {\MARSRevisionCrop{revision/captures/bounded/packet-valid-registers.png}{0.43\linewidth}{0.242187500}{0.313333333}{0.601562500}{0.642222222}};
\node[anchor=north west,inner sep=0,font=\small\bfseries] (btitle) at ([xshift=8mm]atitle.north west -| a.north east) {2. Eight bytes supplied};
\node[anchor=north west,inner sep=0] (b) at (btitle.south west |- a.north west) {\MARSRevisionCrop{revision/captures/bounded/packet-prefix-registers.png}{0.43\linewidth}{0.242187500}{0.313333333}{0.601562500}{0.642222222}};
\foreach \panel in {a,b} {
\begin{scope}[shift={(\panel.south west)},x={($(\panel.south east)-(\panel.south west)$)},y={($(\panel.north west)-(\panel.south west)$)}]
\draw[black,line width=.7pt] (.04,.01) rectangle (.94,.99);
\end{scope}}
\end{tikzpicture}
\caption{Actual return-register observations at \texttt{parse\_ready}, before main stores metadata or exits. The image crops retain native decimal values; the boxes are LaTeX annotations. Availability is a caller-supplied bound even when more bytes are mapped.}
\label{fig:revision-bounded-returns}
\end{figure}


Restore enough available bytes. In disposable copies, introduce both a wrong
magic byte and an unsupported version. Derive which error comes first. Repair
only the magic and repeat, then combine reserved flags with an excessive declared
length. Use the main chapter's priority table to explain each status rather than
accepting whichever comparison happens to run first in your own draft.

For an excessive declaration, set length bytes to hexadecimal \texttt{00 41},
meaning 65. With a complete header, the result is $-5$ even if the declared body
has not arrived. Compare this with the same bytes and available seven, where the
initial result is $-1$. Do not move the length-cap decision after a body read.

Return to the default length three and available eleven. Break at
\texttt{packet\_length} and follow the second availability comparison. The required
count is twelve, so the function returns without entering its checksum loops.
A missing-trailer error must not become an apparent success because padding
happens to contain a suitable checksum byte.

\section{Inspect acceptance before and after publication}

Learn's accepted route uses the unchanged default \texttt{bounded-packet.asm}.
Move one line breakpoint at a time and resume to each requested stop. Learn's
source highlight identifies the next requested stop. If it moves before you
resume, use the selected stack frame in Debug to locate the actual pause. At
\texttt{parse\_packet}, the input, available twelve and selected record are
ready before the first function instruction. At \texttt{packet\_length}, magic,
version and flags have passed; \texttt{\$v0} still holds the tentative flags
error $-4$, and the length field is about to be read. At
\texttt{packet\_fields}, the decoded length is three and the required total
twelve. Only now has the routine established the complete input bound before
sequence decoding and checksum scanning.

The first \texttt{payload\_sum} stop has header total 232, payload total zero,
three bytes remaining and cursor $input+8$. At its next visit, the first byte,
127, has raised those totals to 359 and 127; the cursor moved one byte and two
payload bytes remain. At \texttt{trailer\_sum}, header plus payload totals
742, raw payload sum is 510 and the remaining count is zero. The cursor now
names the trailer. These pauses describe the bytes already visited, before
the selected next load executes.

Continue to \texttt{packet\_commit}. At this point
the checksum residue is zero, but the selected record has not received a store.
Inspect tentative flags, sequence, length and payload sum in the source's temporary
registers. The raw payload sum is 510; its low-byte residue will be 254. The
\texttt{\$v0} register can still hold $-6$ because success has not yet replaced
the tentative checksum status. This is an internal pause, not the return interface.

Step through the successful stores and return to \texttt{parse\_ready}. Inspect
output words $[1,4660,3,254]$, consumed count twelve and payload pointer $input+8$.
Locate the native bytes of the sequence output word. They are
\texttt{34 12 00 00}, while the input sequence field remains \texttt{12 34}.
Use the wire format and native memory representation to explain both observations.

For Learn's checksum-failure route, use \texttt{worked/packet-checksum.asm}.
Only the trailer changes, from hexadecimal \texttt{1A} to \texttt{1B}. Break at
\texttt{trailer\_sum} to compare the same tentative totals, then resume to
\texttt{parse\_ready}. The returned status is $-6$, the payload pointer is zero and the record remains
at its sentinel values. That picture is consistent with no output writes, but
does not prove their absence: the grader checks the store history too. A failed
routine that writes and restores the old words is still incorrect.

\begin{figure}[htbp]
  \centering
  \begin{tikzpicture}[alt={Two actual light-theme Memory observations at parse_ready, each showing rows 0x10010100 and 0x10010110. Native boxes identify the selected sixteen-byte output beginning at 0x10010108. With the valid checksum the words are 1, 4660, 3 and 254. With only trailer 1A changed to 1B, all sixteen selected output bytes remain FF. Surrounding padding and guard bytes are outside the boxes.}]
    \node[anchor=north west,inner sep=0,font=\small\bfseries] (atitle) at (0,0) {1. Accepted packet};
    \node[anchor=north west,inner sep=0] (a) at ([yshift=-2mm]atitle.south west) {\MARSRevisionCrop{revision/captures/bounded/packet-valid-record.png}{0.96\linewidth}{0.032812500}{0.314444444}{0.580468750}{0.608888889}};
    \node[anchor=north west,inner sep=0,font=\small\bfseries] (btitle) at ([yshift=-4mm]a.south west) {2. Incorrect checksum};
    \node[anchor=north west,inner sep=0] (b) at ([yshift=-2mm]btitle.south west) {\MARSRevisionCrop{revision/captures/bounded/packet-checksum-record.png}{0.96\linewidth}{0.032812500}{0.314444444}{0.580468750}{0.608888889}};
    \foreach \panel in {a,b} {
        \begin{scope}[shift={(\panel.south west)},x={($(\panel.south east)-(\panel.south west)$)},y={($(\panel.north west)-(\panel.south west)$)}]
          \draw[black,line width=.7pt] (.592,.319) rectangle (.988,.623);
          \draw[black,line width=.7pt] (.196,.014) rectangle (.592,.319);
        \end{scope}}
  \end{tikzpicture}
  \caption{The selected record spans two native Memory rows. Boxes identify flags/sequence in the first row and length/payload residue in the second. Other bytes include padding and guards. The failed record's unchanged appearance is only a final-state observation; the no-store rule also requires execution-history checks.}
  \label{fig:revision-bounded-records}
\end{figure}


The payload pointer borrows the caller's input storage. In a disposable experiment,
change a payload byte after success and inspect through the returned address.
Explain why the pointer sees the changed byte: no separate payload copy was made.
The caller must preserve the input until it has finished using the slice.

Finish or stop each Debug execution before starting the next worked source.
Learn can present its next question while the earlier execution remains paused;
meeting a checkpoint does not finish that suspended program. When ordinary Run
is requested, all three worked sources exit silently and normally. Their normal
process exit code is independent of the parser's status. Inspect retained
\texttt{status}, \texttt{payload\_address} and the records after Run: accepted
status is twelve with the borrowed pointer, while the two failed parses retain
$-1$ or $-6$ and a zero pointer. Main's metadata stores occur after the earlier
function-return observation.

An edit to a worked or independent source invalidates execution credit tied to
its previous revision. Continue returns to the first activity that needs fresh
evidence while retaining the prediction history. Use disposable copies for
extra faults so their altered behavior does not replace the guided sources.

\section{Diagnose byte and checksum mistakes}

Use a sequence with hexadecimal bytes \texttt{FF 80}. Derive the expected unsigned
value 65408. In a copy, replace the low sequence byte's \texttt{lbu} with
\texttt{lb}. Inspect the sign-extended value before the OR and locate the bits
that corrupt the decoded field. Restore unsigned field decoding and repeat with
an input offset of one byte.

Compare that fault with a payload-only signed-load variant. Each signed high-bit
payload byte differs by minus 256, so both final checksum residues can remain
correct. Explain the arithmetic before deciding whether the variant is valid.
The checks accept equivalent behavior rather than requiring a particular mnemonic.

Introduce an ordinary zero payload byte and recompute the trailer. A parser that
stops at zero will confuse the rest of the body and trailer. Repair the remaining-
count loop instead of adding a text-string terminator rule to this binary format.
Finally exchange the default first two payload values, 127 and 128. Their sum and
the packet checksum stay unchanged; explain what this additive check cannot detect.

Run both original public checks against the complete program:
\begin{verbatim}
java -jar /path/to/cli.jar grade \
  packages/013-binary-packet-parser/checks/parser.kts \
  /path/to/your/bounded-packet.asm
java -jar /path/to/cli.jar grade \
  packages/013-binary-packet-parser/checks/program.kts \
  /path/to/your/bounded-packet.asm
\end{verbatim}
On failure, retain the exact input window, available count and selected output.
Locate the first out-of-window access, wrong decision or forbidden store. A final
status disagreement may originate earlier than the return instruction.

\section{Implement a bounded delimiter parser independently}

Open \texttt{bounded-line.asm} and implement \texttt{read\_line}. Input pointer,
available count and selected eight-byte output pointer arrive in
\texttt{\$a0}/\texttt{\$a1}/\texttt{\$a2}. The input window is inside a 256-byte
arena; output is separate and word aligned. Accept one LF-delimited binary record,
with at most 32 payload bytes. LF is hexadecimal \texttt{0A}; zero, \texttt{0D}
and every other byte value are ordinary payload.

On success, return payload length plus one in \texttt{\$v0} and the original input
pointer in \texttt{\$v1}. Output words are payload length and payload sum modulo
256, excluding LF. Need-more is $-1$; too-long is $-2$. Every failure returns zero
in \texttt{\$v1} and executes no output stores. Preserve the input, guards, other
record, full saved registers, stack pointer and caller-owned stack bytes. Choose
your own temporary registers and private labels.

Begin with zero available bytes and an immediate LF. The first case returns
need-more without reading; the second accepts a zero-length payload and consumes
one byte. The supplied \texttt{41 00 80 0A 42} example consumes four bytes and
returns length three, sum 193. The final \texttt{42} belongs to a later call.

Now write out the capacity boundary. With 32 non-LF bytes and no next byte,
return need-more. If the next supplied byte is LF, accept length 32 and consumed
count 33. If it is not LF, return too-long. Therefore test that next byte for LF
before rejecting it as a 33rd payload byte. Do not read it until it is available.

Stop reading at the first LF, or after 33 bytes when no LF has appeared. Repeated
reads inside this permitted prefix are allowed, so a two-pass search/sum method
can be equivalent to accumulating while searching. Preserve failure outputs by
holding tentative values until acceptance. Do not use the output record as scratch
storage or inspect a following record speculatively.

\noindent\begin{minipage}{\linewidth}
  Run the public checks from the repository root:
  \begin{verbatim}
java -jar /path/to/cli.jar grade \
  revision/lessons/013-bounded-interfaces/independent/bounded-line.kts \
  /path/to/your/bounded-line.asm
\end{verbatim}
\end{minipage}\par\medskip

The starter deliberately stops until implemented. Learn checks the current
source in three activities. First, 139 prefix, boundary and default-caller cases
exercise incomplete input, exact capacity, input offsets, selected records and
complete caller behavior. Then 256 cases cover every possible first-byte value.
Together they retain all 393 original function inputs, add a direct default-input
function case and retain the complete caller. Each function case calls twice
through different continuations, giving 789 checked returns across the profiles.
Each case has a cumulative 5000-instruction budget. The independent algorithm's
private register allocation and frame layout remain your choice.

The final activity requires silent normal exit of the same source. The caller
case has already observed stored status and the full eight-byte borrowed pointer
before the exit selector replaces \texttt{\$v0}. Returning the correct pointer
but storing only its low word would lose information when input lives above
four gigabytes. Checking the assembled address, rather than a fixed low numeric
constant, permits relocation and exposes that truncation.

Expected results and access histories serve different purposes. Read bounds end
at the first LF or the first 33 supplied bytes, even if later bytes are mapped.
Store checks protect caller data, including padding between public objects; a
write followed by restoration is still a forbidden write. Continuation checks
also distinguish returning through the supplied address from jumping to an earlier
continuation and falling through to the requested stop. Both function profiles
and the normal-exit check must pass for the current source revision.

The command-line script above preserves all 394 original checks. The adjacent
JSON profiles are Learn's bounded data, not executable scripts. Keep each failed
input, available count, window offset and selected record. Derive its permitted
read prefix before changing the scan condition. A larger available window still
does not extend this parser's stopping boundary past LF or byte 33.

In disposable copies, use zero as a terminator, then include LF in the payload
sum. Find a concrete counterexample for each and locate its first wrong decision
or update. Finally, write one tentative output word before proving the delimiter
present, then restore it on failure. Explain why the final record can look
unchanged while the no-store contract still fails. Restore the working source
and rerun all three checks before leaving the exercise. These finite samples
support the result; the bounds and delayed-publication argument must cover the
whole stated input domain.

\section{Companion lab: validate UTF-8 across chunks}
\label{revision-lab:utf8}

The packet parser can leave an incomplete request untouched because the caller
still owns the whole prefix. A stream interface may instead accept that prefix
now and receive later bytes in another call. Its state must remember what the
accepted bytes permit next. Resetting at each call would make correctness depend
on where a file reader happened to divide the input.

This companion develops original project \texttt{015-utf8-scanner} through the
native \texttt{015-utf8-streams} unit. Open its assignment and begin with the
complete \texttt{worked/range.asm} experiment. The second complete source,
\texttt{worked/lower-bound.asm}, changes one sample byte so that you can compare
rejection just below a bound with acceptance at that bound.

Implement \texttt{scan\_utf8} in \texttt{independent/stream.asm}. Its allocations,
chunk caller and \texttt{chunk\_ready} observation marker are complete. Replace the
single deliberate \texttt{break} body while retaining that public interface.
Choose your own temporary registers, private labels and equivalent algorithm.
The three inline assessments check direct chunk transitions, interleaved state
objects and the supplied caller. The archived construction regions are not
required by this assignment.

The prerequisite is Chapter 12's bounded byte access and function reasoning.
The new work is a persistent state machine: define which byte can follow a prefix,
carry that decision across calls, count completed values, and retain the first
failure. This scanner validates and counts UTF-8. It does not produce decoded
text or replacement characters.

\subsection*{Distinguish bytes, scalar values and displayed text}

A Unicode code point is a numbered position, conventionally written with
\texttt{U+} and hexadecimal digits. A Unicode \emph{scalar value} is a code point
outside the surrogate range U+D800 through U+DFFF. UTF-8 represents a scalar with
one through four bytes. The restrictions in Table~\ref{tab:revision-utf8-leads}
exclude surrogate encodings, values beyond U+10FFFF and encodings longer than
needed.\footnote{The Unicode Consortium, \emph{The Unicode Standard}, version
17.0.0, Section 3.9.3 and Table 3-7:
\url{https://www.unicode.org/versions/Unicode17.0.0/core-spec/chapter-3/}.}

A complete scalar is the unit this lab counts. It is not necessarily one visible
character. For example, U+0065 followed by U+0301 represents an \texttt{e} and a
combining acute accent: two scalars that can form one grapheme cluster in displayed
text. U+00E9 is one scalar with a precomposed accented letter. This scanner counts
them as two and one respectively; it does not normalize either representation or
segment grapheme clusters.\footnote{For grapheme clusters and text boundaries,
see Unicode Standard Annex 29, revision 47, Section 3:
\url{https://www.unicode.org/reports/tr29/tr29-47.html}.}

The supplied stream has eleven bytes, written here in hexadecimal:
\[
\underbrace{\mathtt{41}}_{\mathrm{U+0041}}\quad
\underbrace{\mathtt{E2\ 82\ AC}}_{\mathrm{U+20AC}}\quad
\underbrace{\mathtt{F0\ 9F\ 99\ 82}}_{\mathrm{U+1F642}}\quad
\underbrace{\mathtt{00}}_{\mathrm{U+0000}}\quad
\underbrace{\mathtt{C3\ A9}}_{\mathrm{U+00E9}}.
\]
It contains five complete scalars. Byte zero is a valid one-byte encoding; it
neither ends this supplied stream nor suppresses the following scalar. U+FEFF,
unassigned scalar values and noncharacters also count when their encodings are
well formed. Assignment and display policies are outside this validator's contract.
The program exits silently, so inspect its state rather than waiting for printed
characters.

\subsection*{Derive the next-byte rule from an accepted prefix}

A one-byte encoding has high bit zero. A multibyte lead specifies its width;
following bytes, called \emph{continuation bytes}, have bit pattern
\texttt{10xxxxxx}. Their ordinary hexadecimal interval is \texttt{80} through
\texttt{BF}. A continuation cannot begin a new scalar. Bytes \texttt{C0},
\texttt{C1} and \texttt{F5} through \texttt{FF} cannot begin a valid encoding
either. The first continuation after four particular leads has a tighter bound.

\begin{table}[htbp]
  \centering
  \small
  \begin{tabular}{llll}
    \toprule
    Lead bytes (hex) & Width & First continuation (hex) & Later continuations (hex) \\
    \midrule
    00--7F           & 1     & none                     & none                      \\
    C2--DF           & 2     & 80--BF                   & none                      \\
    E0               & 3     & A0--BF                   & 80--BF                    \\
    E1--EC, EE--EF   & 3     & 80--BF                   & 80--BF                    \\
    ED               & 3     & 80--9F                   & 80--BF                    \\
    F0               & 4     & 90--BF                   & 80--BF                    \\
    F1--F3           & 4     & 80--BF                   & 80--BF                    \\
    F4               & 4     & 80--8F                   & 80--BF                    \\
    \bottomrule
  \end{tabular}
  \caption{UTF-8 lead classes and the ranges used by this scanner. Width includes the lead. Only the first continuation can have a restricted range; later continuations use the ordinary range. Derived from Unicode 17.0.0 Table 3-7.}
  \label{tab:revision-utf8-leads}
\end{table}


These bounds do more than check the continuation tag. After \texttt{E0}, accepting
\texttt{80} would allow a value below the three-byte minimum. After \texttt{ED},
accepting \texttt{A0} would enter the surrogate range. The \texttt{F0} lower bound
excludes values that need only three bytes; the \texttt{F4} upper bound excludes
values above the Unicode limit. After a valid first continuation, reset the range
to \texttt{80} through \texttt{BF} for every remaining continuation. The tighter
rule does not persist for the whole sequence.\footnote{The byte grammar is also
specified in RFC 3629, Section 4:
\url{https://www.rfc-editor.org/rfc/rfc3629.html\#section-4}.}

You can derive a boundary numerically. A three-byte encoding carries four useful
bits from the lead and six from each continuation. For bytes $b_0,b_1,b_2$, its
value is
\[
 ((b_0\mathbin{\&}15)\ll12)
 \mathbin{|}((b_1\mathbin{\&}63)\ll6)
 \mathbin{|}(b_2\mathbin{\&}63).
\]
Thus \texttt{E0 A0 80} gives $32\cdot64=2048$, hexadecimal 0800, the first
three-byte scalar. \texttt{E0 9F 80} would give 1984, which has a shorter encoding.
The scanner need not assemble this numerical value: retaining the allowed range
and remaining count is enough to validate it. Use the formula to explain the
range before encoding comparisons in assembly.

Load input as unsigned bytes when comparing these ranges. An \texttt{lb} of
\texttt{E0} sign-extends a negative value; unsigned comparisons with small positive
bounds would then classify the wrong number. Advance the byte pointer with
full-width arithmetic even though this lab's byte counts fit within a word.
Neither a short chunk nor a small scalar implies a small address.

\subsection*{Isolate an inclusive byte comparison}

The complete range experiment loads byte \texttt{9F} and calls this helper with
lower bound \texttt{A0} and upper bound \texttt{BF}. The arguments are ordinary
unsigned byte values in registers, so decimal 159, 160 and 191 describe the same
comparison. The helper returns one for membership and zero otherwise. Its result
is separate from the scanner's status codes.

\noindent\begin{minipage}{\linewidth}
\MARSCodeLabel{Test a byte against inclusive unsigned bounds}
\begin{minted}{mars-idea}
within_range:
    sltu $t0, $a0, $a1
    bnezc $t0, outside_range
    nop
    sltu $t0, $a2, $a0
    xori $v0, $t0, 1
    jrc $ra
outside_range:
    move $v0, $zero
    jrc $ra
\end{minted}
\end{minipage}

The first \texttt{sltu} produces one if the byte is below the lower bound. The
compact branch then selects the zero-return path. If that test is zero, the second
comparison asks whether the upper bound is below the byte. That comparison is one
only for a byte that is too large; \texttt{xori} reverses its zero-or-one result.
Equality with either bound therefore remains accepted.

Use the first guided Debug sequence to observe \texttt{9F}: the lower-bound test
is one and the returned result is zero. In \texttt{worked/lower-bound.asm}, the
sample becomes \texttt{A0}. Its lower-bound test is zero and the returned result
is one. Both callers store the answer in \texttt{range\_result} and exit silently.
The public \texttt{worked/range.kts} script checks every byte value against seven
relevant intervals as well as the first complete caller.

A correct interval helper is only one part of a scanner. The accepted prefix
chooses the interval, the pending count determines whether a scalar is complete,
and the call's final flag decides what an unfinished prefix means at the end.
Those decisions need persistent state, which the helper deliberately does not
provide.

\subsection*{Give persistent state a precise meaning}

The function takes the selected state pointer in \texttt{\$a0}, current chunk
pointer in \texttt{\$a1}, byte count in \texttt{\$a2}, and final flag in
\texttt{\$a3}. The flag is zero or one. A chunk contains zero through 64 bytes
inside the supplied 96-byte arena and can begin at any byte offset that leaves
room for it. State is a separate, word-aligned 32-byte object, either
\texttt{state\_a} or \texttt{state\_b}. Calls receive a reachable state; you do
not need to accept arbitrary corrupted descriptors. For an active call, old total
plus supplied count is at most 4096.

\begin{table}[htbp]
\centering
\small
\begin{tabular}{rll}
\toprule
Byte offset & State word & Meaning at return\\
\midrule
0 & total & Accepted bytes over all calls\\
4 & count & Completed scalars over all calls\\
8 & need & Continuations still needed\\
12 & minimum & Lowest permitted next continuation\\
16 & maximum & Highest permitted next continuation\\
20 & start & Pending or failed lead's absolute offset; $-1$ if idle\\
24 & status & 0 active, 1 finished, $-1$ malformed, $-2$ truncated\\
28 & error & Detecting offset, or $-1$ before any failure\\
\bottomrule
\end{tabular}
\caption{The selected scanner state consists of eight native words. Total counts accepted bytes, so it also names the offset where the next byte would be accepted. A terminal rejection returns $-3$ without replacing the stored status.}
\label{tab:revision-utf8-state}
\end{table}


The empty state is $[0,0,0,128,191,-1,0,-1]$, in the table's word order and decimal
notation. An idle active state has need zero, ordinary continuation bounds and
start $-1$. A pending active state has need from one through three and retains
the absolute offset of its lead. Its bounds describe the \emph{next} byte, not
the first byte of the next call. Error remains $-1$ while active.

At each input-byte boundary, total is the number of accepted bytes over all calls;
count is the number of complete scalars among them. A valid pending prefix
contributes to total but not yet to count. The current call's accepted count is
separate. Return it in \texttt{\$v1}; total includes earlier calls too.

If need is zero, an ASCII byte completes a scalar immediately. A legal multibyte
lead records start at the old total, establishes need and the first-continuation
range, and is then accepted. If need is positive, test the current byte against
the stored range. An accepted continuation reduces need and restores ordinary
bounds. Reaching need zero completes one scalar and resets start to $-1$.
Only acceptance advances total, the call's accepted count and the input pointer.
An impossible byte advances none of them.

Keep this invariant at the loop's top. First determine whether the current call
has an unexamined byte, then load within its supplied interval. A zero-length
call performs no input load. The contract permits read-ahead within the current
chunk, but never across its end, even if the next chunk is mapped immediately
after it. A chunk ending with \texttt{E2} does not grant access to the following
\texttt{82}; retain the pending rule for the next call instead.

\subsection*{Trace one scalar across two chunk boundaries}

The caller supplies the eleven-byte example as four chunks, with the last one
marked final. Predict every row of Table~\ref{tab:revision-utf8-chunks} before
running your implementation. All offsets are zero based; state and return values
are decimal unless the byte column says otherwise.

\begin{table}[htbp]
  \centering
  \small
  \setlength{\tabcolsep}{5pt}
  \begin{tabular}{rlrrrrll}
    \toprule
    Index & Bytes (hex) & Total & Count & Need & Start & Next range & $(v_0,v_1)$ \\
    \midrule
    0     & 41 E2       & 2     & 1     & 2    & 1     & 128--191   & $(0,2)$     \\
    1     & 82 AC F0    & 5     & 2     & 3    & 4     & 144--191   & $(0,3)$     \\
    2     & 9F 99       & 7     & 2     & 1    & 4     & 128--191   & $(0,2)$     \\
    3     & 82 00 C3 A9 & 11    & 5     & 0    & $-1$  & 128--191   & $(1,4)$     \\
    \bottomrule
  \end{tabular}
  \caption{Completed calls for the supplied four-chunk stream. Only index three is final. Stored status equals $v_0$ in these active calls; error stays $-1$. The returned accepted count is local to each call, while total and count accumulate.}
  \label{tab:revision-utf8-chunks}
\end{table}


The first chunk accepts \texttt{A} and the \texttt{E2} lead. Total becomes two,
but only one scalar is complete. The second chunk finishes U+20AC and accepts
\texttt{F0}, so need becomes three and the next-byte minimum becomes 144,
hexadecimal \texttt{90}. The third chunk accepts two continuations. The first,
\texttt{9F}, satisfies that tighter bound; after it, minimum returns to 128.
Need is still one when the chunk ends. Its two bytes have not completed another
scalar.

The fourth chunk begins with \texttt{82}, which completes U+1F642. It then accepts
the zero byte and \texttt{C3 A9}. Count becomes five, need becomes zero and the
final flag permits successful completion. Splitting the same bytes at another
position must produce the same final total and scalar count, although intermediate
states and per-call accepted counts will differ.

Break at \texttt{chunk\_ready} in Debug. The caller has already stored
\texttt{chunk\_index}, \texttt{last\_status} and \texttt{last\_used}; their values
identify the completed call. The state words must agree with that same call.
For the native Debug activity, place the breakpoint on the \texttt{nop} beneath
\texttt{chunk\_ready} and set its condition to \texttt{\$s3 == 0}. After that
stop, change only the condition to \texttt{\$s3 == 1}, then resume. Repeat for
2 and 3. Keep the same execution session and source throughout: restarting would
return to the first chunk, while editing source would invalidate the observation. The revised caller uses
compact transfers and preserves its live values in saved registers. Your scanner
must preserve those registers and the stack pointer across all returns.

Figure~\ref{fig:revision-utf8-pending} compares the pending rules after chunks
one and two in actual native Memory observations. Read the boxed words as need
and minimum using the state layout above.

\begin{figure}[htbp]
\centering
\begin{tikzpicture}[alt={Actual Memory observations at chunk_ready after chunks one and two. The first state begins with words total five, count two, need three and minimum 144. The next begins with total seven, count two, need one and minimum 128. Both retain maximum 191, lead start four, active status zero and error minus one. Native boxes identify need and minimum.}]
\node[anchor=north west,inner sep=0,font=\small\bfseries] (atitle) at (0,0) {1. After chunk 1: a four-byte lead accepted};
\node[anchor=north west,inner sep=0] (a) at ([yshift=-2mm]atitle.south west) {\MARSRevisionCrop{revision/captures/utf8/utf8-four-byte-lead.png}{0.96\linewidth}{0.032812500}{0.314444444}{0.580468750}{0.608888889}};
\node[anchor=north west,inner sep=0,font=\small\bfseries] (btitle) at ([yshift=-4mm]a.south west) {2. After chunk 2: one continuation still needed};
\node[anchor=north west,inner sep=0] (b) at ([yshift=-2mm]btitle.south west) {\MARSRevisionCrop{revision/captures/utf8/utf8-pending-tail.png}{0.96\linewidth}{0.032812500}{0.314444444}{0.580468750}{0.608888889}};
\foreach \panel in {a,b} {
\begin{scope}[shift={(\panel.south west)},x={($(\panel.south east)-(\panel.south west)$)},y={($(\panel.north west)-(\panel.south west)$)}]
\draw[black,line width=.7pt] (.592,.319) rectangle (.988,.623);
\end{scope}}
\end{tikzpicture}
\caption{Actual light-theme states for the supplied stream. Boxes mark need and next minimum: three/144 becomes one/128, while count stays two. The first allowed continuation resets the tighter minimum. Each four-byte group is one native little-endian word.}
\label{fig:revision-utf8-pending}
\end{figure}


\subsection*{Separate malformed input, truncation and terminal reuse}

A nonfinal chunk that ends in a still-possible prefix returns zero. A final chunk
with no pending prefix returns one; a final chunk with a pending prefix returns
$-2$, meaning truncated. The latter is a promise that no more bytes will arrive,
not a new byte to validate. Even an empty final call can therefore finish or
truncate a stream.

An impossible byte returns $-1$ immediately, whether or not that call is final.
Store its absolute offset in error, retaining the lead offset in start. For an
invalid byte encountered while idle, that byte itself supplies start. Accepted
progress preceding the error stays committed. This differs from the packet
parser's no-output-on-failure contract: the stream's state records both its
accepted prefix and its terminal diagnostic.

Compare final input \texttt{41 E0} with final input \texttt{41 E0 80}. Both
accept two bytes and complete one scalar. The first ends with a possible prefix:
status $-2$, start one, error two. In the second, byte \texttt{80} at offset two
is already impossible after \texttt{E0}: status $-1$, start one, error two. Do not
call it truncated merely because the chunk was final. The offending byte is not
included in \texttt{\$v1} or total. A final flag never replaces an earlier byte
error. Figure~\ref{fig:revision-utf8-errors} compares those two completed calls.

\begin{figure}[htbp]
\centering
\begin{tikzpicture}[alt={Actual Memory observations for final chunks 41 E0 80 and 41 E0. Both states have total two, scalar count one, need two, minimum 160, maximum 191, start one and detecting offset two. The malformed state's status word is FFFFFFFF, minus one; the truncated state's is FFFFFFFE, minus two. Native boxes mark status and detecting offset.}]
\node[anchor=north west,inner sep=0,font=\small\bfseries] (atitle) at (0,0) {1. Final bytes 41 E0 80: malformed};
\node[anchor=north west,inner sep=0] (a) at ([yshift=-2mm]atitle.south west) {\MARSRevisionCrop{revision/captures/utf8/utf8-malformed-final.png}{0.96\linewidth}{0.032812500}{0.314444444}{0.580468750}{0.608888889}};
\node[anchor=north west,inner sep=0,font=\small\bfseries] (btitle) at ([yshift=-4mm]a.south west) {2. Final bytes 41 E0: truncated};
\node[anchor=north west,inner sep=0] (b) at ([yshift=-2mm]btitle.south west) {\MARSRevisionCrop{revision/captures/utf8/utf8-truncated-final.png}{0.96\linewidth}{0.032812500}{0.314444444}{0.580468750}{0.608888889}};
\foreach \panel in {a,b} {
\begin{scope}[shift={(\panel.south west)},x={($(\panel.south east)-(\panel.south west)$)},y={($(\panel.north west)-(\panel.south west)$)}]
\draw[black,line width=.7pt] (.592,.014) rectangle (.988,.319);
\end{scope}}
\end{tikzpicture}
\caption{The same accepted prefix has different terminal causes. Native boxes mark status and detecting offset. In the first picture, \texttt{FF FF FF FF} stores $-1$; in the second, \texttt{FE FF FF FF} stores $-2$. Both detect at offset two, but only the first has an impossible byte there. These are completed calls at \texttt{chunk\_ready}.}
\label{fig:revision-utf8-errors}
\end{figure}


A successful final call, a malformed call and a truncated final call all make
the state terminal. On any later call, return $(-3,0)$ without reading input or
writing caller data. Preserve the stored status and error, so an attempted reuse
does not erase the original result. Check status before entering the active scan.
A zero-length later call has the same rejection rule. To begin an unrelated stream,
the caller must deliberately initialize a fresh empty state.

While active, stores may change only the selected state. Preserve the input,
guards, other state, caller metadata and full saved registers. A correctly restored
private stack frame is permitted. Active store order and temporary registers are
not prescribed. A terminal rejection permits no caller-data stores, including a
write that restores the same value. Keep terminal rejection outside any common
active publication path that would rewrite state unconditionally.

\subsection*{Build tests around the first changed decision}

Begin with ASCII and an empty nonfinal call, then one multibyte scalar at each
width. Divide each valid encoding after every possible byte and insert empty
nonfinal chunks. Next test the four restricted leads on both sides of their first-
continuation bounds. Include \texttt{E0 A0 80} and \texttt{F0 90 80 80} to expose
a routine that forgets to restore ordinary bounds after the first continuation.

Test \texttt{ED A0} and \texttt{F4 90} before supplying any further byte: both
are already malformed. Test lone \texttt{80}, \texttt{C0} and \texttt{F5} while
idle. Repeat a pending prefix once with final zero and once with final one, then
try another call after each terminal outcome. Keep the original failure status
and offsets in view during that last experiment.

Use two state objects for interleaved independent streams. Feed \texttt{E2} to
one and \texttt{C3} to the other; then finish them with \texttt{82 AC} and
\texttt{A9} respectively. Their saved pending rules belong to those state objects,
not to private global variables or a register value left over from the last call.
Choose a fresh input offset for the second pair of calls. The supplied checks also
change scratch registers between calls, so register residue cannot serve as
persistent state.

Use \emph{Check this program} in the native lesson for the 426 original stream
cases, twelve additional interleaved-state cases and fifteen caller variants.
The interleaving assessment alternates both state objects within one case, checks
their complete state after each call, changes temporary registers and revisits
terminal outcomes. All three assessments use the same continuing source.

The two original public check domains are also available through the CLI. From
the downloaded unit's directory, run these commands on your completed source:

\noindent\begin{minipage}{\linewidth}
\begin{verbatim}
java -jar /path/to/cli.jar grade \
  independent/stream.kts independent/stream.asm
java -jar /path/to/cli.jar grade \
  independent/caller.kts independent/stream.asm
\end{verbatim}
\end{minipage}

The additional interleaving cases run through the native assessment. Finish with
an ordinary Run of the four-chunk caller; its successful silent exit is separate
from the state and access checks above.


When a check fails, retain the stream prefix, selected state, current chunk and
final flag. Find the first byte whose permitted range or accepted count diverges.
Then compare the saved state before and after that call. A correct final scalar
count cannot excuse an out-of-chunk load, an overwritten earlier error or lost
pending state. Keep one boundary trace and one malformed-versus-truncated trace
with your passing results, explaining why changing chunk boundaries does not
change the validity of the complete stream.

